\documentclass[11pt,a4paper]{article}
\usepackage[margin=2.5cm]{geometry}
\usepackage[T1]{fontenc}
\usepackage{lmodern}
\usepackage{parskip}
\pagestyle{empty}

\begin{document}

\begin{center}
{\large\bfseries EH-605: Modeling Earth System and Sustainability\\[2pt] Project Outline}\\[12pt]
{\Large\bfseries Why Is Methane Rising Again? A Methane--OH Box Model of Recent Trends and the Global Methane Pledge}\\[10pt]
Shardul Junagade \quad (Roll No. 23110297)
\end{center}

\vspace{6pt}

\textbf{Project idea.}
Atmospheric methane, the second largest contributor to greenhouse warming after CO$_2$, almost stopped rising between 2000 and 2006, started growing again in 2007 and reached record growth rates in 2020 and 2021. Emission inventories do not show the same pattern. This project will use simple box models from the course to understand what controls these changes.

\textbf{Objectives.}
\begin{enumerate}
\item Test whether a one-box model with a constant lifetime can reproduce the observed methane record since 1984.
\item Build a coupled methane--OH model and study its feedback, stability and the possibility of a tipping point.
\item Estimate how much the Global Methane Pledge could lower methane concentration and warming by 2050.
\end{enumerate}

\textbf{Approach.}
The models will be written as ordinary differential equations, solved in Python and fitted to NOAA observations using EDGAR emission data. The fitted model will then be used to compare future emission scenarios and to look for a cost-effective way to reduce emissions.

\textbf{Expected outcomes.}
A model that explains the observed growth rate, an estimate of the methane--OH feedback strength, and a comparison of the pledge with an economically optimal path, along with a short look at emissions from India.

\end{document}
